\documentclass[]{report}

\usepackage{import}

\import{../}{settings.tex}

\title{Fonctions holomorphes}
\author{Cours de LETELLIER Emmanuel, écrit par SADR TAHOURI Luca}
\date{}

\begin{document}

    \maketitle
    \tableofcontents
    
\chapter{Fonctions holomorphes}

\section{Fonctions Holomorphes}

\subsection{Définitions}

Soit U un ouvert de $\C$\\

\begin{df} 
    Soit $f:U \to \C$. On dit que f est $\C$-dérivable en $z_0$ si il existe $\alpha \in \C$ tel que
\[\lim\limits_{h \to 0} \left(\frac{f(z_0+h)-f(z_0)}{h}=\alpha\right) (\iff \forall \epsilon >0, \exists\eta>0,|h|<\eta \implies|g(h)-\alpha|<\epsilon)\]
\end{df}

\begin{df} 
    Si une fonction $f:U \to \mathbb{C}$ est holomorphe s'il est $\mathbb{C}$-dérivable en tout point de U. On note:
$\begin{array}{ccccc}
f' & : & E & \to & F \\
 & & z & \mapsto & f'(z) \\
\end{array}$ la fonction dérivée
\end{df} 

\begin{tm} 
    Si $f:U \to \mathbb{C}$ est holomorphe, alors $f'$ est continue
\end{tm}

\begin{proof}
    admis
\end{proof}

\begin{re}{}{}
    $\lim\limits_{h \to 0} \left(\frac{f(z_0+h)-f(z_0)}{h}\right)=\alpha \iff f(z_0+h)=f(z_0)+\alpha h +(h)\epsilon (h)$ avec $\epsilon(h)\underset{h\to 0}{\longrightarrow} 0$. On voit alors que si $f$ est holomorphe, alors $f$ est continue.
\end{re}

\begin{prop}La fonction $\begin{array}{ccccc}
f' & : & \mathbb{C} & \to & \mathbb{C} \\
 & & z & \mapsto & \bar{z} \\
\end{array}$ n'est $\mathbb{C}$-dérivable en aucun point
\end{prop}

\begin{proof} $\lim\limits_{h \to 0} \left(\frac{f(z_0+h)-f(z_0)}{h}\right)=\lim\limits_{h \to 0} \frac{\bar{h}}{h}$\\
Or $\lim\limits_{h \to 0,h\in \mathbb{R}} \frac{\bar{h}}{h}=1,\lim\limits_{h \to 0,h\in \mathbb{C}} \frac{\bar{h}}{h}=-1 $, il n'y a donc pas de limite.
\end{proof}

\begin{re}On identifie $\begin{array}{ccccc}
\iota & : & \mathbb{C} & \to & \mathbb{C} \\
 & & a+ib & \mapsto & (a,b) \\
\end{array}$. Dans ce cas, l'ouvert U de $\mathbb{C}$ correspond à un ouvert $\widetilde{U}$ de $\mathbb{R}^2$ et on note $\begin{array}{ccccc}\widetilde{f} & : & \widetilde{U} & \to & \mathbb{R}^2\end{array}$ la fonction qui correspond à $\begin{array}{ccccc}f & : & U & \to & \mathbb{C}\end{array}$. La condition $f(z_0+h)=f(z_0)+\alpha h+ |h|\epsilon(h)$ avec $z_0=x_0+iy_0$, $h=s+it$ et $\alpha = a+ib$ devient:\\
$\widetilde{f}(x_0+s,y_0+t)=\widetilde{f}(x_0,y_0)+\begin{pmatrix}
a & -b \\
b & a 
\end{pmatrix}\begin{pmatrix}
s \\
t
\end{pmatrix}+|(s,t)|\widetilde{\epsilon}(s,t)$. Donc si f est $\mathbb{C}$-dérivable en $z_0$ alors $\widetilde{f}$ est différentiable en $(x_0,y_0)$:\\
$\begin{array}{ccccc}
d\widetilde{f} & : & \mathbb{R}^2 & \to & \mathbb{R}^2 \\
 & & \begin{pmatrix}
u \\
v
\end{pmatrix} & \mapsto & \begin{pmatrix}
a & -b \\
b & a 
\end{pmatrix}\begin{pmatrix}
u \\
v
\end{pmatrix} \\
\end{array}$. Réciproquement si $\widetilde{f}$ est $\mathbb{R}$-différentiable en $(x_0,y_0)$ et que sa différentielle $\begin{array}{ccccc}d_{(x_0,y_0)}\widetilde{f} & : & \mathbb{R}^2 & \to & \mathbb{R}^2\end{array}$ est $\mathbb{C}$-linéaire, alors $f$ est $\mathbb{C}$-dérivable en $z_0=x_0+iy_0$. En effet, si $d\widetilde{f}$ est $\mathbb{C}$-linéaire, sa matrice dans la base canonique est de la forme $\begin{pmatrix}
a & -b\\
b & a
\end{pmatrix}$
\end{re}

\begin{re}[Calcul] $\begin{array}{ccccc}
\widetilde{\varphi} & : & \R^2 & \to & \R^2 \\
 & & \begin{pmatrix}
x \\
y
\end{pmatrix} & \mapsto & \begin{pmatrix}
a & b \\
c & d 
\end{pmatrix}\begin{pmatrix}
x \\
y 
\end{pmatrix} \\
\end{array}$, $\begin{array}{ccccc}
\varphi & : & \C & \to & \C \\
 & & x+iy & \mapsto & ax+by+i(cx+dy) \\
\end{array}$\\
\begin{align*}
\varphi(i(x+iy)) & =\varphi(-y+ix)\\ 
&=i(ax+by)-(cx+dy) \\
&=(bx-ay)+i(dx-cy) \\
& \iff \begin{cases} ax+by=-cy+dx & \forall x,y\\ cx+dy=ay-bx \end{cases} \\
& \implies \begin{cases} \text{Si } x=0 \text{ et } y=1  & \implies b=-c \\ \text{Si } x=1 \text{ et } y=0  & \implies a=d  \end{cases}
\end{align*}
\end{re}

\begin{prop}[Proposition (Équations de Cauchy-Riemann)]
    Soit $\begin{array}{ccccc}f & : & U & \to & \C\end{array}$ et $\begin{array}{ccccc}
\widetilde{f} & : & \widetilde{U} & \to & \R^2 \\
 & & (x,y) & \mapsto & (P(x,y),Q(x,y)) \\ 
\end{array}$ \\
Alors $f$ est $\C$-dérivable en $z_0=x_0+iy_0 \in U$ si et seulement si P et Q sont $\R$-différentiable en $(x_0,y_0)$ et $\begin{cases} \frac{\delta P}{\delta x}(x_0,y_0)=\frac{\delta Q}{\delta y}(x_0,y_0)\\ \frac{\delta P}{\delta y}(x_0,y_0)=-\frac{\delta Q}{\delta x}(x_0,y_0) \end{cases}$
\end{prop}

\begin{proof}
    Si $\widetilde{f}$ est différentiable en $z_0=(x_0,y_0)$. La matrice jacobienne de $\widetilde{f}$ en $(x_0,y_0)$\\
    \[\begin{pmatrix}
\frac{\delta P}{\delta x}(x_0,y_0) & \frac{\delta P}{\delta y}(x_0,y_0) \\
\frac{\delta Q}{\delta x}(x_0,y_0) & \frac{\delta Q}{\delta y}(x_0,y_0)
    \end{pmatrix}\]\\
\end{proof}

\begin{re}
    Pour vérifier que $\begin{array}{ccccc}
f : & U & \to & \C \\
  & x+iy & \mapsto & P (x, y) + iQ (x, y)  
\end{array}$ avec $\begin{array}{ccccc}P,Q  : & \widetilde{U} & \to & \R\end{array}$ est dérivable en $z_0=x_0+iy_0$, on vérifie les équations de Cauchy Riemann.
\end{re}

\begin{ex}
    $\fundfto{f}{\C}{\C}{x+iy}{x-iy}$. On a $P(x,y)=x$, $Q(x,y)=-y$\\
On a donc $\begin{cases} \frac{\delta P}{\delta y}(x,y)=\frac{\delta Q}{\delta x}(x,y)=0\\ \frac{\delta P}{\delta x}=1 \neq -1 = \frac{\delta Q}{\delta y} \end{cases}$ donc f n'est pas différentiable.
\end{ex}

\subsection{Quelques propriétés}

Soit U un ouvert de $\C$

\begin{df}
    $\mathcal{H}(U)=\{$fonctions holomorphes de $U \to \C \}$ est une algèbre unitaire.
\end{df}

\begin{prop}
    $\mathcal{H}(U)$ est un $\C$-espace vectoriel et on a les propriétés suivantes:\\
    \begin{itemize}
        \item $(f+g)'=f'+g'$ et $(\lambda f)'=\lambda f'$
        \item Si $f,g \in \mathcal{H}(U)$ alors $fg\in \mathcal{H}(U)$ et $(fg)'=f'g+fg'$
        \item Si $f\in \mathcal{H}(U)$ et $f$ ne s'annule pas, alors $\frac{1}{f} \in \mathcal{H}(U)$
        \item Si $\begin{array}{ccccc} f : & U & \to & V \end{array}$ et $\begin{array}{ccccc} g : & V & \to & \C \end{array}$ holomorphes, alors $\begin{array}{ccccc} g \circ f : & U & \to & \C \end{array}$ est holomorphe et $(g \circ f)'=f'\times(g' \circ f)$
    \end{itemize}
\end{prop}

\begin{ex}
    Fonctions polynomiales $\mfun{z}{\sum_{i=1}^{k} a_i z^i}$ est holomorphe sur $\C$. La fonction $\begin{array}{ccccc}z & \mapsto  \frac{1}{z}\end{array}$ est holomorphe sur $\C^*$. Plus généralement, une fraction $\frac{P(z)}{Q(z)}$ avec $P,Q\in \C[X]$ est holomorphe en dehors des zéros de Q.
\end{ex}

\begin{prop}
    Soit $U\subseteq \C$ un ouvert connexe. $f\in \mathcal{H}(U)$ est constante si et seulement si $f'=0$.
\end{prop}

\begin{proof}
    $\implies$ Immédiate \\
    $\impliedby$ supp($f'$)$=0$. Soient $\funto{P,Q}{\widetilde{U}}{\R}$ les parties réelles et imaginaires telles que $f(x+iy)=P(x,y)+iQ(x,y)$. Montrons que P et Q sont constantes. Comme f est holomorphe, P et Q sont différentiables et $f'(z_0)=\frac{\delta P}{\delta x}(x_0,y_0)-i\frac{\delta P}{\delta y}(x_0,y_0)$.\\
    Comme $f'$ est continue, les dérivées $\frac{\delta P}{\delta x}$, $\frac{\delta P}{\delta y}$, $\frac{\delta Q}{\delta x}$ et $\frac{\delta Q}{\delta y}$ sont continues.\\ 
    Donc P et Q sont de classe $C^1$. Or, d'après le théorème des accroissements finis, $\forall(a,b),(a',b')\in U, \exists(c,d)\text{ sur }\left. \right]A,A'\left[ \right.$ tel que\\
    $P(a,b)-P(a',b')=\frac{\delta P}{\delta x}(c,d)(a'-a)+\frac{\delta P}{\delta y}(c,d)(b'-b)$.
    Comme $f'=0$, on a $\frac{\delta P}{\delta x}=\frac{\delta P}{\delta y}=0$ donc $P(a,b)=P(a',b')$ donc $P(a,b)=P(a',b')$ donc P constante (de même pour Q)
\end{proof}

\section{Séries entières}

Soit $(a_n)_{n\in\N}$ suite de nombres complexes et $\sum_{n\ge0}a_nz^n$ la série formelle associée.

\begin{df}
    Le rayon de convergence de $\sum_{n\ge0}a_n z^n$ est $R=\sup \{r\ge0,|a_n|r^n $ bornée$\}$
\end{df}

\begin{tm}
    Soit R rayon de convergence de $\sum_{n\ge0}a_n z^n$:
    \begin{enumerate}
        \item $\forall z \in \C $ tel que $|z|<R$, la série $\sum_{n\ge0}a_nz^n$ est absolument convergente
        \item $\forall z\in\C$ tel que $|z|>R$, la série $\sum_{n\ge0}a_nz^n$ diverge
    \end{enumerate}
\end{tm}

\begin{prop}[Proposition (Règle de d'Alembert)]
    $\sup (a_n \ne) 0$ pour $n \gg0$. Si $|\frac{a_{n+1}}{a_n}| \underset{n\to +\infty}{\longrightarrow}l \in \R^+\cup\{+\infty\}$ alors le rayon de convergence est $R=\frac{1}{l}$ (avec les conventions $R=0$ si $l=+\infty$, $R=+\infty$ si $l=0$)
\end{prop}

\begin{prop}[Proposition (Règle de Cauchy)]
    Si $\lim\limits_{n \rightarrow +\infty}|a_n|^{\frac{1}{n}}=l \in [0,+\infty]$. Alors le rayon de convergence de $\sum_{n\ge0}a_nz^n$ est $R=\frac{1}{l}$. Elle se généralise de la manière suivante :\\
    limite supérieur et limite inférieur : Soit $(U_n)_{n\in\N}$ suite à valeurs dans $\R$ ou $\bar{\R}=[-\infty,+\infty]$, $v_n=\sup \{U_k,k\ge n\}$ et $w_n=\inf \{U_k, k\ge n\}$. On a $(v_n) \searrow$ et $(w_n) \nearrow$. Elles admettent une limite dans $\bar{\R}$
\end{prop}

\begin{prop}
    Supposons que $\sum_{n\ge0}a_nz^n$ a un rayon de convergence $R>0$ et considérons la fonction $f(z)=\sum_{n\ge0}a_nz^n$ pour $|z|<R$, alors :
    \begin{enumerate}
        \item Le rayon de convergence de la série dérivée $\sum_{n\ge1}na_nz^{n-1}$ est R
        \item La fonction $f$ est holomorphe sur $D(0,R)=\{z\in\C, |z|<R\}$ et pour $z \in D(0,R)$, on a $f'(z)=\sum_{n\ge1}na_nz^{n-1}$
    \end{enumerate}
\end{prop}

\begin{proof}
    \begin{enumerate}
        \item Soit $R'$ rayon de convergence de $\sum_{n\ge1}na_nz^{n-1}$. Supposons $R<R'$, donc $\exists z_0$ tq $R<|z_0|<R'$. La série $\sum_{n\ge1}na_nz_0^{n-1}$ est absolument convergente et $|na_nz^{n}|\le n|a_nz_0^{n-1}|\times |z_0|$ d'où la convergence absolue de $\sum_{n\ge0}a_nz_0^{n}$ ce qui contredit $|z|>R'$.\\
        Supposons que $R'<R$. $\exists l\in\R $ tel que $R'<l<R$ et $z_0\in\C$, $R'<|z_0|<l<R$, comme $l<R$, la suite $a_nl^n$ bornée, donc $\exists M$ tel que $\forall n \in \N,|a_n|l^n \le M$ d'où $n|a_nz_0^n|=|a_n|l^n(\frac{|z_0|}{l})^n\le nMO^n$ avec $O=\frac{|z|}{l}<1$. Le terme $nMO^n$ est celui d'une série convergente, donc $\sum_{n\ge1}n|a_n|z_0^{n}$ est convergente d'où $\sum_{n\ge1}n|a_n|z_0^{n-1}$ convergente.
        \item Soit $z\in D(0,R)$. Montrons que f est $\C$-dérivable en z calculons $f'(z)$. On choisit $r$ tel que $|z|<r<R$. Pour $h\in\C$ tq $|z|+|h|<r$, on a $\frac{f(z+h)-f(z)}{h}-\sum na_nz^{n-1}=\frac{\sum a_n (z+h)^n-\sum a_nz^n}{h}-\sum na_nz^{n-1}=\sum_{n\ge1}a_n(\frac{(z+h)^n-z^n}{h}-nz^{n-1})=\sum_{n\ge2}a_nv_n(h)$. Pour n fixé, $v_n(h) \underset{h\to 0}{\longrightarrow} 0 $(Laisser en exercice). On veut montrer que $\lim\limits_{h \rightarrow 0} \sum a_nv_n(h)=0$. Or $v_n(h)=(z+h)^{n-1}+(z+h)^{n-2}+\ldots+z^{n-1}-nz^{n-1}$. On utilise la formule $x^n-y^n=(x-y)(x^{n-1}+\ldots+y^{n-1})$. On a donc pour h tel que $|z|+|h|<r$, $a_nv_n(h)\le zn|a_n|r^{n-1}$ terme d'une série convergente. Donc $\sum|a_nv_n(h)|$ converge normalement
    \end{enumerate}
\end{proof}

\begin{re}
    $f'(z)=\sum_{n\ge1}na_nz^{n-1}$\\
    $f''(z)=\sum_{n\ge2}n(n-1)a_nz^{n-2}$\\
    \vdots \\
    $a_n=\frac{f^{(n)}(0)}{n!}$ pour tout n. On note :\\
    $\sum_{n\ge0}\frac{f^{(n)}(0)}{n!}z^n$ série de Taylor de $f$ en 0.
\end{re}

\begin{cor}
    Si $\begin{array}{ccccc} f : & D(0,R) & \to  \C\end{array}$ est somme de $\sum a_n z^n$, alors cette série est la série de Taylor de $f$ en 0.
\end{cor}

\section{Fonctions analytiques}

\begin{df}
    U ouvert de $\C$ et $\begin{array}{ccccc} f : & U & \to  \C\end{array}$. On dit que $f$ est analytique si pour $z_0 \in U$, il existe $D(z_0,r)\subseteq U$ et une série entière de rayon de convergence $R\ge r$ tel que $\forall z \in D(z_0,r),f(z)=\sum_{n\ge 0}a_n(z-z_0)^n$
\end{df}

\begin{prop}
    Soit $\begin{array}{ccccc} f : & U & \to  \C\end{array}$ une fonction analytique, alors f est holomorphe et admet des dérivées complexes de tout ordre qui sont aussi des fonctions holomorphes.
\end{prop}

\begin{re}
    Les coefficients $a_n$ dépendent de $z_0$, et $a_n=\frac{f^{(n)}(z_0)}{n!}$
\end{re}

\begin{proof}[Démonstration de la proposition]
    Si $\begin{array}{ccccc} f : & U & \to  \C\end{array}$ est analytique et $z_0\in U,\forall z \in D(z_0,r),f(z)=\sum a_n(z-z_0)^n$. Pour $h\in D(0,r),g(h)=\sum_{n\ge 0}a_nh^n$ est holomorphe et $a_n=\frac{g^{(n)(0)}}{n!}$\\
    $D(z_0,r){\longrightarrow} D(0,r)$ par $z\mapsto z-z_0$ et $\begin{cases}
        D(z_0,r)\underset{f}{\longrightarrow}\C\\
        D(0,r)\underset{g}{\longrightarrow}\C
    \end{cases}$ Donc f est holomorphe et $\frac{f^{(n)}(z_0)}{n!}=a_n=\frac{g^{(n)}(z)}{n!}$
\end{proof}

\begin{prop}
    Soient $\sum a_nz^n$ série entière de rayon de convergence $R>0$ et $\begin{array}{ccccc}
f : & D(0,R) & \to & \C \\
  & z & \mapsto & \sum_{n\ge0}a_nz^n  
\end{array}$ alors f est analytique $\iff \forall z_0 \in D(0,R)$, 
$\exists z>0$ tel que $D(z_0,r) \subseteq D(0,R)$ et \\$f(z)=\sum_{n\ge0}\frac{f^{(n)}(0)}{n!}(z-z_0)^n$
\end{prop}

\begin{proof}
    Soit $z_0\in D(0,R)$\\
    \begin{align*}f^{(p)}(z_0)&=\sum_{n=p}^{+\infty}n(n-1)\ldots(n-p+1)a_nz_0^{n-p}\\
    f^{(p)}(z_0)&=\sum_{q=0}^{+\infty}\frac{(p+q)!}{q!}a_{p+q}z_0^q\text{ avec le changement de variable $n-p=q$}\\
    |f^{(p)}(z_0)|&\leq \sum_{q=0}^{+\infty}\frac{(p+q)!}{q!}|a_{p+q}||z_0^q|
    \end{align*}
\end{proof}

\begin{ra}
    Soit $I$ dénombrable et $(U_i)_{i\in I} \subset R^+$. On dit que $\sum U_i$ est sommable si $\{\sum_{i\in F}U_i |F_{fini}\subseteq I\}$ est borné.
\end{ra}

\begin{prop}
    $I=\bigsqcup_{n\in \N}I_n$, alors $(U_i)_{i\in I}$ est sommable \ssi $\forall n\in\N, (U_i)_{i\in I_n}$ est sommable et $\sum_n(\sum_{i\in I_n}U_i)=\sum_{i\in I }U_i:=\sup\{\sum_{i\in F}U_i|F_{fini}\subseteq I\}$
\end{prop}

\begin{ex}
    $I=\N^2,(U_i)_{i\in I}$, qu'en est-t'il de $\sum U_i $ ?\\
    \begin{enumerate} 
        \item 
            \[\text{$(U_i)_{i\in I}$ sommable $\iff$}\cho{\forall n \sum_{m=0}^{+\infty}U_{(n,m)} \text{ convergente}\\
            \sum_{n\ge0}(\sum_{m\ge0} U_{(n,m)}) \text{ converge}}\]
        \item $I_n'=\{(p,q)\in I \text{ tels que } p+q=n\}$\\
            $I=\bigsqcup_{n\ge0}I_n',(U_i)_{i\in I} $ sommable$\iff\forall n,\sum_{(p,q)\in I_n'}$ converge et $\sum_{n\ge0}\sum_{p+q=n}U_{(p,q)}$ converge
    \end{enumerate}
\end{ex}

\begin{proof}[Suite de la démonstration]
    Soit $r<R-|z_0|$, $\sum_{p,q=0}^{+\infty}\frac{(p+q)!}{p!q!}|a_{p+q}||z_0^q|r^p=^?\sum_{p+q=n}\frac{n!}{p!q!}|z_0^q|r^p=(|z_0|+r)^n$\\
    Comme $r<R-|z_0|$. $\sum a_n(z_0+r)^n$ pour $z_0\in D(0,R)$ absolument convergente donc convergente.\\
    Donc la série de gauche converge et \\
    $\sum_{p=0}^{\infty}\frac{1}{p!}|f^{(p)}(z_0)|r^p\le\sum_{q=0}^{\infty}\sum_{p=0}^{\infty}\frac{(p+q)!}{p!q!}|a_{p+q}||z_0^q|r^p$\\
    $\sum_{p=0}^{\infty}\frac{1}{p!}|f^{(p)}(z_0)|r^p\le \sum_{q=0}^{+\infty}\frac{(p+q)!}{q!}|a_{p+q}||z_0^q|$\\
    Donc la série de Taylor de $f$ en $z_0$ est convergente sur $D(z_0,R-|z_0|)$ $(r=|z-z_0|)$\\
    Question : $f(z)=\sum_{p\ge0}f^{(p)}(z_0)(z-z_0)^p,\forall z\in D(z_0,R-|z_0|)\iff \forall h \in D(0,R-|z_0|),f(z_0+h)=\sum_{p\ge 0}\frac{f^{(p)}(z_0)}{p!}f^p$. Soit $|h|<R-|z_0|$. On a vu que $S(h)=\sum_{p,q}\frac{(p+q)!}{p!q!}a_{p+q}z_0^qh^p$ est alors absolument convergente.\\
    \begin{align*} 
S(h) & = \sum_{p}\frac{1}{p!}(\sum_{q}a_{p+q}\frac{(p+q)!}{q!}z_0^q)h^p\\
& =\sum_{p} \frac{1}{p!}f^{(p)}(z_0)h^p 
\end{align*}\\
     \begin{align*}
S(h) & = \sum_{n\ge0}a_n\sum_{p+q=n}\frac{n!}{p!q!}z_0^qh^p \\
& = \sum_n a_n (z_0+h)^n \\
&=f(z_0+h) & \text{d'où l'égalité}
\end{align*} 
\end{proof}

\chapter{Exponentielle et logarithme complexe}

\section{Fonction exponentielle}

\begin{df}
    Une fonction entière est une fonction holomorphe sur $\C$
\end{df}

\begin{prop}
    La série entière $\sum_{n=0}^{+\infty}\frac{z^n}{n!}$ définit une fonction entière $\begin{array}{ccccc}
\exp : & \C & \to & \C \\
  & z & \mapsto & \exp (z)=e^z   
\end{array}$\\
\begin{enumerate}
    \item On a $e^0=1$ et $\exp'=\exp$
    \item $\forall z, \omega \in \C, e^{\bar{z}}=\bar{e^z} $ et  $e^{z+\omega}=e^ze^\omega$ et donc $\begin{array}{ccccc}
\exp : & \C & \to & \C^*
\end{array}$
    \item La restriction de $\exp$ à $\R$ est l'exponentielle réelle qui est un difféomorphisme $\begin{array}{ccccc}
\R & \to &  \R_+^*
\end{array}$ et on a $|e^z|=1$ \ssi $z\in i\R$
\end{enumerate}
\end{prop}

\begin{proof}
    \begin{enumerate}
        \item Posons $a_n=\frac{1}{n!}$. Comme $\frac{a_{n+1}}{a_n} \longrightarrow 0$ donc d'après d'Alembert, $\sum_{n\ge0}\frac{z^n}{n!}$ converge pour tout $z\in\C$. Il est clair que $\exp(0)$.\\
        $\exp'$ est donnée par la série dérivée qui est aussi $\sum\frac{z^n}{n!}$ et donc $\exp'=\exp$.
        \item $S_N(z)=\sum_{n=0}^N\frac{z^n}{n!}$\\
        $S_N(z)=\sum_{n=0}^N\frac{z^n}{n!}$ comme $z \mapsto \bar{z}$ est continue, on a $\overline{\lim\limits_{N\to+\infty}S_n(z)}=\lim\limits_{N\to \infty}S_N(\overline{z})$ donc $\overline{e^z}=e^{\bar{z}}$.\\
        Montrons que $e^{z+\omega}=e^ze^\omega$ :\\
        Pour $a,b\in \C,g(z)=e^ze^{a+b-z}$, $g'(z)=e^ze^{a+b-z}-e^ze^{a+b-z}=0$ donc g constante, \\$g(z)=g(0)=e^{a+b}$ donc $e^ze^{a+b}e^{-z}=e^{a+b}$\\
        Si $z=a,e^ae^b=e^{a+b}$ donc $e^a\neq0, \forall a\in\C, e^{a+ib}=e^ae^{ib}$\\
        Si $|e^{a+ib}|=1$ alors $e^a|e^{ib}|=1\implies a=0$\\
        $|e^{ib}|^2=e^{ib}\overline{e^{ib}}=e^{ib}e^{-ib}=1$ donc $|e^{ib}|=1$ donc $z\in i\R\implies |z|=1$
    \end{enumerate}
\end{proof}

\begin{re}
    $\exp$ est l'unique application entière $\begin{array}{ccccc}
f : & \C & \to & \C 
\end{array}$ qui vérifie :$\begin{cases}
    f(0)=1\\
    f'=f
\end{cases}$\\
En effet, si $\begin{array}{ccccc}
g : & \C & \to & \C \\
  & z & \mapsto & f(z)e^{-z}  
\end{array}$ alors $g'=0$ car ($g'=f'(z)e^{-z}-f(z)e^{-z}=0$) donc g constante, donc $g(z)=g(0)=f(0)e^{-0}=1,\forall z \in \C$, donc $f(z)=e^z$
\end{re}

Nous voulons montrer que $\begin{array}{ccccc}
\exp : & \C & \to & \C^*
\end{array}$ est surjective

\begin{prop}[Proposition (Inversion locale holomorphe)]
    Soit $\begin{array}{ccccc}
f : & U & \to & \C
\end{array}
$ holomorphe et $z_0\in U$ tels que $f'(z_0)\neq 0$. Alors $f$ est biholomorphisme au voisinage de $z_0$, c'est à dire qu'il existe un voisinage $V\subseteq U$ de $z_0$\tq $f(V)$ soit ouvert et $\begin{array}{ccccc}
V & \to & f(V)
\end{array}$ bijective d'inverse $\begin{array}{ccccc}
f^{-1} : & f(V) & \to & V
\end{array}$ est holomorphe.
\end{prop}

\begin{proof}
    Soit $\begin{array}{ccccc}
\widetilde{f} : & \widetilde{U} & \to & \R^2
\end{array}$ fonction correspondante à $f, f'(z_0)=a+ib\neq 0$.\\
La matrice jacobienne de $\widetilde{f}$ en $(a,b)$ est $\begin{pmatrix}
a & -b \\
b & a
\end{pmatrix}$ qui est inversible donc d'après le théorème d'inversion locale, $\exists \widetilde{V}$ ouvert de $\R^2$ et $\widetilde{W}$ ouvert de $\widetilde{U}$ \tq $(a,b)\in \widetilde{W}\underset{\widetilde{f}}{\longrightarrow} \widetilde{V}$ et la différentielle de $\widetilde{f}$ en $(a,b)$ est un isomorphisme comme elle est $\C$-linéaire, son inverse aussi.
\end{proof}

\begin{tm}
    $\begin{array}{ccccc}
\exp : & (\C,+) & \to & (\C^*,\times)
\end{array}$ est un morphisme de groupe surjectif. Il existe un unique réel positif, noté $\pi$ \tq $\ker(\exp)=2\pi i \Z$.\\
Tout nombre complexe non nul a une décomposition polaire $z=re^{i\theta}$ avec $|z|=r$\\
\begin{itemize}
    \item$e^{i\theta}=\cos(\theta)+i\sin(\theta)$
    \item $e^{i(\theta+\theta')}=e^{i\theta}e^{i\theta'}$
\end{itemize}
\end{tm}

\begin{proof}
    On a déjà vu que $\exp$ est un morphisme de groupes et que c'est une application ouverte.\\
    Surjectivité : $\Im(\exp):=H$ est un sous groupe ouvert de $\C^*$ et \[\C^*=H\bigsqcup(\bigcup_{\substack{U \text{ ouvert}\\U \subset A}}aH)\] \\
Chaque classe $aH$ est aussi un ouvert de $\C^*$ (car $\begin{array}{ccccc}
& \C^* & \to & \C^* \\
& z & \mapsto & az  
\end{array}$ est un homéomorphisme). Comme $\C^*$ est connexe, on a que $V=\varnothing $($\iff \C^*=H$) et donc $\exp$ est surjective. Noyau de $\exp$ : On a vu que $\ker(\exp)\subseteq i \R$. Il s'agit de trouver le noyau de $\begin{array}{ccccc}
h : & (\R,+) & \to & (\S^1,*) \\
  & t & \mapsto & e^{it}  
\end{array}$ $h$ est continu donc $\ker h$ est un sous groupe fermé de $\R$. (Rappel : Les sous groupes additifs de $\R$ sont les $\alpha \Z , \alpha\in\R^+$ et les sous groupes denses dans $\R$) Donc $\ker(h)=\R$ ou $\ker h=\alpha\Z$ pour $\alpha\in\R^+$. Comme h est surjective, $\ker h \neq\R$ donc $\ker h = \alpha \Z, \alpha\in \R^+$, $\ker h \neq 0$ ($e^{it}=ie^{4it}=1$) donc $\alpha>0$, et on définie $\pi$ pour $\alpha=2\pi$.
\end{proof}

\begin{cor}
    Si $\begin{array}{ccccc}
f : & U & \to & \C
\end{array}$ est une application holomorphe dont la dérivée ne s'annule jamais, alors $f$ est une application ouverte.
\end{cor}

\begin{proof}
    Il faut montrer que $U'$ ouvert $\implies f(U')$ ouvert.\\
    Soit $U'$ ouvert de $U$ et $y\in f(U')$ qu'on écrit $y=f(x),x\in U'$ comme $f'(x)\neq0$, il existe un voisinage $W$ de x ($W \subseteq U'$) tel que $\begin{array}{ccccc}
& W & \to & f(W)
\end{array}$. Il existe y et f(W) ouvert dans $\C$ qui est contenu dans $f(U')$ et qui contient y donc $f(U')$ ouvert de $\C$
\end{proof}

\section{Logarithme}

$\begin{array}{ccccc}
\exp : & \R & \to & \R^*
\end{array}$ est une bijection croissante, son inverse est le logarithme néperien noté :\\
$\begin{array}{ccccc}
\ln : & \R^*_+ & \to & \R
\end{array}$. Dans le cas complexe, $\begin{array}{ccccc}
\exp : & \C & \to & \C^*
\end{array}$ n'est pas injective. Si $z\in \C^*,z=e^w=e^{a+ib}=e^ae^{ib},a=\log |z|$

\begin{df}
    $U\subseteq\C^*$ ouvert.\\
    Une application $\begin{array}{ccccc}
f : & U & \to & \C
\end{array}$ est une détermination du logarithme lorsque :\begin{enumerate}
    \item f est continue
    \item $\forall z\in U,e^{f(z)}=z$
\end{enumerate}
\end{df}

\begin{prop}
    Il n'existe pas de détermination du logarithme sur $\C^*$ tout entier.
\end{prop}

\begin{proof}
    Si on en avait une, cela voudrait dire qu'on a une détermination continue.\\ Soit $\begin{array}{ccccc}
\theta : & \C^* & \to & \R \\
  & z & \mapsto & \theta(z)  
\end{array}$ et $\begin{array}{ccccc}
f : & \C^* & \to & \C \\
  & z & \mapsto & \log(z)+i\theta(z)  
\end{array}$\\
En particulier si $z=e^{it}$, on aurait $e^{it}=e^{i\theta(e^{it})}$.
On pose $\begin{array}{ccccc}
\delta : & \R & \to & \R \\
  & t & \mapsto & t-\theta(e^{it})  
\end{array}$. Dans ce cas $\Im(\delta)\subseteq2\pi\Z$, comme $\delta$ est continue et $2\pi\Z$ discret, donc $\delta$ est constante. $\forall t, t-\theta(e^{it})=a\iff \theta(e^{it})=t-a$ (or $\theta(e^{it})$ est périodique mais $t-a$ ne l'est pas)
\end{proof}. 

\begin{prop}
    Soient $U\subseteq \C^*$ ouvert connexe et $\begin{array}{ccccc}
f_0 : & U & \to & \C \\ 
\end{array}$ une détermination du logarithme. Les autres déterminations du logarithme sur $U$ sont données par $f_n=f_0+2in\pi,n\in\Z$
\end{prop}

\begin{proof}
    Exercice
\end{proof}

\begin{prop}
    Soient $U\subseteq\C$ ouvert connexe. \begin{enumerate}
        \item Si $\begin{array}{ccccc}
f : & U & \to & \C
\end{array}$ est une détermination du $\log$ sur U, alors f est holomorphe \\et on a $f'(z)=\frac{1}{z},\forall z\in U$.
    \item Si U connexe. Soit $\begin{array}{ccccc}
g : & U & \to & \C
\end{array}$ holomorphe tel que $g'(z)=\frac{1}{z}$ pour tout $z\in U$. Alors il existe une constante $\alpha\in\C$ tel que $\begin{array}{ccccc}
& U & \to & \C \\
& z & \mapsto & g(z)-\alpha  
\end{array}$ soit une détermination du log
    \end{enumerate}
\end{prop}

\begin{proof}
    Voir p.13
\end{proof}

\begin{df}
    Soit $\begin{array}{ccccc}
f : & U & \to & \C
\end{array}$ une fonction continue. Une primitive $F$ de $f$ sur $U$ est une fonction holomorphe $\begin{array}{ccccc}
F : & U & \to & \C
\end{array}$ telle que $F'=f$
\end{df}

\begin{re}
    La fonction $\begin{array}{ccccc}
& \C^* & \to & \C^* \\
& z & \mapsto & \frac{1}{z}
\end{array}$ n'admet pas de primitive
\end{re}

\section*{Détermination principale du logarithme}

    $\begin{array}{ccccc}
\exp : & \C & \to & \C^* \\
\end{array}$ est un morphisme de groupes surjectif de noyau $2i\pi\Z$. Sa restriction $\begin{array}{ccccc}
\exp : & \left\{ \omega\in\C,-\pi\le\Im(\omega)<\pi\right\} & \to & \C^*\\
\end{array}$? L'image de la droite $\left\{\omega\in\C,\Im(\omega=-\pi)\right\}$ est la demi-droite $\R^*\subseteq \C$.\\
$\begin{array}{ccccc}
\exp : & \left\{ \omega\in\C,-\pi<\Im(\omega)<\pi\right\} & \to & \C^*\backslash\R^*_- \\
\end{array}$ est holomorphe bijective dont la dérivée est partout non nulle. Elle réalise donc un biholomorphisme entre deux ouverts.\\
L'application réciproque : $\begin{array}{ccccc}
\log_p : & \C\backslash\R^+ & \to & \left\{\omega\in\C,-\pi<\Im(z)<\pi\right\} \\
\end{array}$ s'appelle la détermination principale du logarithme.

\begin{lm}
    $\log_p$ est développable en série entière sur le disque ouvert $D(1,1)$ centré en 1 et de rayon 1.$$\log_p(z)=\sum_{n=0}^{+\infty}(-1)^n\frac{(z-1)^{n+1}}{n+1}$$
\end{lm}

\begin{proof}
    Voir p.15
\end{proof}

\section{Fonctions trigonométriques et hyperboliques}

Pour $z\in\C$, \begin{align*}
    \sin(z)=\frac{e^{iz}-e^{-iz}}{2i} && \cos(z)=\frac{e^{iz}+e^{-iz}}{2}\\
    \sinh(z)=\frac{e^{z}-e^{-z}}{2} && \cosh(z)=\frac{e^{z}+e^{-z}}{2}
\end{align*}

\chapter{Intégration complexe}

\section{Arcs paramétrés}

\subsection{Définitions}

On appelle arc paramétré, ou simplement arc, une application continue $\begin{array}{ccccc}
\gamma : & [t_0,t_1] & \to & \C \\ 
\end{array}$, $(t_0,t_1)\in\R$.\\
L'origine de $\gamma$ est le point $z_0=\gamma(t_0)$ et l'extremité $z_1=\gamma(t_1)$.\\
Si $z_0=z_1$, on dit que $\gamma$ est un lacet.\\
Le support de $\gamma$ est $\left\{\gamma(t),t\in[t_0,t_1]\right\}=\Im(\gamma)$ et on le notera supp$(\gamma)$.

\subsection{Exemples}

\begin{enumerate}
    \item Soient $a,b\in\C$, le segment $[a,b]$ est le support de l'arc $\begin{array}{ccccc}
\gamma : & [0,1] & \to & \C \\
  & t & \mapsto & (1-t)a+tb  
\end{array}$.\\
Il partage le même support avec $\begin{array}{ccccc}
\gamma' : & [0,1] & \to & \C \\
  & t & \mapsto & ta+(1-t)b
\end{array}$
    \item Ligne polygonale : On peut aligner plusieurs segments vu précédemment de cette manière $$\gamma(t)=\left\{\begin{array}{ll}
 (1-t)z_0+tz_1 & \text{ si } 0 \le t \le 1\\
 (2-t)z_1+tz_2 & \text{ si } 1 \le t \le 2\\
 (3-t)z_2+tz_3 & \text{ si } 2 \le t \le 3
 \end{array}\right.$$
    \item Le cercle $x^2+y^2=R^2$ :\\
    C'est le support de $\begin{array}{ccccc}
\gamma : & [0,2\pi] & \to & \C \\
  & t & \mapsto & Re^{it}  
\end{array}$ ou $\begin{array}{ccccc}
\gamma : & [0,2\pi] & \to & \C \\
  & t & \mapsto & Re^{-it}
\end{array}$
\end{enumerate}

\subsection{Opération sur les arcs}

\begin{df}
    Etant donné un arc $\begin{array}{ccccc}
\gamma : & [t_0,t_1] & \to & \C
\end{array}$. \\L'arc opposé $-\gamma$ est défini par $\begin{array}{ccccc}
-\gamma : & [t_0,t_1] & \to & \C \\
  & t & \mapsto & \gamma(t_0+t_1-t)
\end{array}$(qui commence à $t_1$ et qui finit à $t_0$)\\
Deux arcs $\gamma$ et $\gamma'$ définis sur $[t_0,t_1]$ et $[t_0',t_1']$ sont consécutifs si $\left\{\begin{array}{cc}
t_1=t_0'\\
\gamma(t_1)=\gamma'(t_0')
\end{array}\right.$\\
Dans ce cas, on peut définir $\begin{array}{ccccc}
\gamma+\gamma' : & [t_0,t_1] & \to & \C \\
& t & \mapsto & \left\{\begin{array}{cc}
\gamma(t)     & \text{ si } t_0 \le t \le t_1 \\
\gamma'(t)     & \text{ si } t_0' \le t \le t_1'
\end{array}\right.  
\end{array}$
\end{df}

\subsection{Chemins et circuits}

Un chemin est un arc $\gamma$ qui peut s'écrire sous la forme $\gamma=\gamma_1+\dots+\gamma_n$ où chaque $\gamma_i$ est continûment dérivable (de classe $C^1$).\\
Un circuit est un lacet qui est un chemin.

\subsection{Longueur}

    Etant donné un arc $\begin{array}{ccccc}
\gamma : & [t_0,t_1] & \to & \C
\end{array}$. A toute subdivision $t_0=x_0<\dots<x_n=t_1$, on peut associer la ligne polygonale de sommets $z_k=\gamma(x_k)$\\
La longueur de la ligne polygonale est $|z_1-z_0|+|z_2-z_1|+\dots|z_{k-1}-z_k|$.\\
La longueur de $\gamma$ est la borne supérieure de toutes les longueurs de ses lignes polygonales. On la note $L(\gamma)$.\\
On dit que $\gamma$ est rectifiable si $L(\gamma)<+\infty$.

\begin{prop}
    Si $\gamma$ est un chemin défini sur $[t_0,t_1]$, alors $\gamma$ est rectifiable et $$L(\gamma)=\int_{t_0}^{t_1} |\gamma'(t)| \,dt$$
\end{prop}

\section{Intégrale le long d'un chemin}

\subsection{Définitions}

Soit $\begin{array}{ccccc}
\gamma : & [t_0,t_1] & \to & \C
\end{array}$ un chemin et $\begin{array}{ccccc}
F : & \text{supp}(\gamma) & \to & \C
\end{array}$ une fonction continue. On définie alors $$\int_{\gamma} F(z) \,dz=\int_{t_0}^{t_1} F(\gamma(t))\gamma'(t) \,dt$$

\subsection{Exemples}

\begin{enumerate}
    \item Si $\gamma(t)=t,\forall t\in [t_0,t_1]$, alors $$\int_{\gamma} F(z) \,dz=\int_{t_0}^{t_1} F(t) \,dt$$
    \item Si $\gamma(t)=Re^{it},\forall t\in[0,2\pi]$, alors $$\int_{\gamma}\frac{dz}{z}=\int_{0}^{2\pi} \frac{Rie^{it}}{Re^{it}} \,dt=i\int_{0}^{2\pi} dx=2\pi i$$
\end{enumerate}

\subsection{Propriétés}

\begin{itemize}
    \item $\int_{\gamma} (F_1(z) +F_2(z))\,dz=\int_{\gamma} F_1(z) \,dz+\int_{\gamma} F_2(z) \,dz$
    \item $\int_{\gamma} \lambda F(z) \,dz=\lambda\int_{\gamma} F(z) \,dz$
    \item $\int_{\gamma_1+\gamma_2} F(z) \,dz=\int_{\gamma_1} F(z) \,dz+\int_{\gamma_2} F(z) \,dz$
    \item $\int_{-\gamma} F(z) \,dz=-\int_{\gamma} F(z) \,dz$
\end{itemize}

\begin{lm}
    $$\left|\int_{\gamma} F(z) \,dz\right|\le L(\gamma)\underset{z\in\text{supp}(\gamma)}{\sup}(|F(z)|)$$
\end{lm}

\begin{proof}
    Voir p.18
\end{proof}

\subsection{Changement de paramétrage}

Soit $\begin{array}{ccccc}
\gamma : & [t_0,t_1] & \to & \C
\end{array}$ un chemin, et soit $\begin{array}{ccccc}
\varphi : & [u_0,u_1] & \to & [t_0,t_1]
\end{array}$ surjective de classe $C^1$ telle que $\varphi(y_0)=t_0$ et $\varphi(u_1)=t_1$. \\
Alors $\delta:=\gamma\circ\varphi$ est un chemin, et on dit que $\delta$ est un reparamétrage de $\gamma$.\\
On a : $$\int_\gamma F(z)dz=\int_\delta F(z)dz$$
C'est une conséquence de la formule de changement de variable.
\begin{align*}
\int_{t_0}^{t_1} f(t) \,dt = & \int_{u_0}^{u_1} f(\varphi(u))\varphi'(u)\,du\\
\int_\gamma F(z)\,dz= & \int_{t_0}^{t_1} F(\gamma(t))\gamma'(t)\,dt \\
\int_{\delta} F(z) \,dz=& \int_{u_0}^{u_1} F(\delta(u))\delta'(u)\,du\\
=& \int_{u_0}^{u_1} F(\gamma\circ\varphi(u))\varphi'(u)(\gamma'\circ\varphi)(u)\,du
\end{align*}
Appliquer le changement de variable à $f(t)=F(\gamma(t))\gamma'(t)$

\begin{align*}
     \int_{\delta} F(z)\,dz & = \int_{u_0}^{u_1} f(\varphi(u))\varphi'(u)\,du\\
    & =\int_{t_0}^{t_1}f(t)\,dt=\int_\gamma F(z)\,dz
\end{align*}

\section{Expression intégrale de fonctions analytiques}

\subsection{Passage à la limite}

\begin{lm}
    On considère un chemin $\gamma$ et une suite $(F_n)_{n\in\N}$ de fonctions complexes "continues sur le support de $\gamma$" et $(F_n)$ converge uniformément vers $F$, alors $$\int_\gamma F(z)\,dz=\lim_{n\to+\infty}\int_\gamma F_n(z)\,dz$$
\end{lm}

\begin{proof}
    Voir p.19
\end{proof}

\subsection{Une formule intégrale}

\begin{prop}
    Étant donné un chemin $\gamma$ et une fonction complexe $\varphi$ continue sur supp$(\gamma)$, on définit : $$\begin{array}{ccccc}
F : & \C\backslash\text{supp}(\gamma) & \to & \C \\
  & z & \mapsto & \int_\gamma \frac{\varphi(u)}{u-z} \,du
\end{array}$$
\begin{center}
    C'est une fonction analytique, et de plus\\
    $$F^{(n)}(z)=n!\int_\gamma\frac{\varphi(u)}{u-z}\,du$$
\end{center}
\end{prop}

\begin{proof}
    Voir p.20
\end{proof}

\chapter{Primitives des fonctions complexes}

On n'admet plus le théorème : $\begin{array}{ccccc}
f : & U & \to & \C
\end{array}\text{ holomorphe} \implies \begin{array}{ccccc}
f' : & U & \to & \C
\end{array} \text{ continue}$

\section{Primitives}

Soit U un ouvert de $\C$ et $\begin{array}{ccccc}
f : & U & \to & \C
\end{array}$ holomorphe. Rappelons qu'une primitive de f sur U est une fonction holomorphe telle que sa dérivée est f.

\begin{ex}
    \begin{enumerate}
        \item Le polynôme $P(z)=a_nz^n+\dots+a_1z+a_0$ admet sur $\C$ des primitives de la forme : $\frac{a_n}{n+1}z^{n+1}+\frac{a_{n-1}}{n}z^n+\dots+a_0z+C,C\in\C$
        \item Sur $\Omega=\C\backslash\R^-$ la détermination principale du logarithme est une primitive de $\frac{1}{z}$.
    \end{enumerate}
\end{ex}

\begin{prop}
    Soit $R$ le rayon de convergence de $\sum a_nz^n$. Soit $$\begin{array}{ccccc}
f : & D(0,R) & \to & \C \\
  & z & \mapsto &  \displaystyle{\sum_0^{+\infty}a_nz^n} \end{array}$$
Alors $$\begin{array}{ccccc}
F : & D(0,R) & \to & \C \\
  & z & \mapsto & \displaystyle{\sum_0^{+\infty}\frac{a_n}{n+1}z^{n+1}}  
\end{array}\text{ est une primitive de f}$$
\end{prop}

\begin{proof}
    Voir p.21
\end{proof}

\begin{cor}
    Toute fonction analytique admet localement des primitives.
\end{cor}

\section*{Primitives et intégration}

\begin{prop}
    Soient $\begin{array}{ccccc}
f : & U & \to & \C
\end{array}$ continue et $\begin{array}{ccccc}
F : & U & \to & \C
\end{array}$ une primitive de $f$. Etant donné un chemin $\gamma$ tracé dans $U$ d'origine $z_0$ et d'extrémité $z_1$.$$\int_\gamma f(z)\,dz=F(z_1)-F(z_0)$$
\end{prop}

\begin{proof}
    Voir p.22
\end{proof}

\begin{cor}
    Si la fonction $\begin{array}{ccccc}
f : & U & \to & \C
\end{array}$ continue admet une primitive sur $U$, on a pour tout circuit fermé $\gamma$ dans $U$, $$\int_\gamma f(z)\,dz=0$$
\end{cor}

\begin{ex}
    $\begin{array}{ccccc}
& \C^* & \to & \C \\
& z & \mapsto & \frac{1}{z}  
\end{array}$ continue. On a vu qu'en posant $\gamma(t)=e^{it},t\in[0,2\pi]$, on a $$\int_\gamma\frac{dz}{z}=2\pi i$$
Donc $\frac{1}{z}$ n'a pas de primitive sur $\C^*$.
\end{ex}

\section{Primitives sur un ouvert étoilé}

\subsection{Ensemble étoilé}

\begin{df}
$E\subseteq\C,E\neq\emptyset$. On appelle centre de $E$ tout point $e_0\in E$ qui vérifie la condition :$$\forall z\in E, [z,z_0]\subseteq E$$
On dit alors que $E$ est étoilé par rapport à $z_0$. Un ensemble est étoilé s'il a au moins un centre
\end{df}

\begin{ex}
    Un ensemble est convexe si chacun de ses points est un centre.
\end{ex}

\subsection{Primitives sur un ouvert étoilé}

Soit $\Omega$ un ouvert étoilé et $z_0$ un centre de $\Omega$. Soit $\begin{array}{ccccc}
f : & \Omega & \to & \C
\end{array}$

\begin{lm}
    Deux primitives de $f$ sur $\Omega$ diffèrent d'une constante.
\end{lm}

\begin{re}
    On avait montré en première partie du cours cet énoncé (ouvert connexe) en admettant que la dérivée d'une fonction holomorphe est continue (ce qu'on ne suppose plus dans cette deuxième partie)
\end{re}

\begin{proof}
    Voir p.23
\end{proof}

\begin{prop}
    Pour que $f$ admette des primitives sur $\Omega$, il suffit qu'on ait $\displaystyle{\int_\theta f(z)\,dz=0}$ pour tout circuit triangulaire $\theta$ tracé dans $\Omega$.
\end{prop}

\begin{re}
    La réciproque a déjà été vu : Si $f$ admet une primitive, alors $\displaystyle{\int_\theta f(z)\,dz=0}$ pour tout circuit $\theta$ dans $\Omega$.
\end{re}

\begin{proof}
    Voir p.24
\end{proof}

\begin{cor}
    Si $\begin{array}{ccccc}
f : & \Delta & \to & \C
\end{array}$ est continue avec $\Delta$ un ouvert de $\C$ et si $\displaystyle{\int_\theta}f(z)\,dz=0$ pour tout circuit triangulaire $\theta$ tracé dans $\Delta$, alors $f$ admet localement des primitives.
\end{cor}

\subsection{Primitives sur un ouvert connexe (domaine)}

\begin{ra}
    Un ouvert connexe de $\C$ est connexe par arcs
\end{ra}

\begin{ex}
    $\C^*$ est un ouvert connexe mais ce n'est pas un ouvert étoilé.
\end{ex}

\begin{lm}
    Si $\begin{array}{ccccc}
f : & \Delta & \to & \C
\end{array}$ est continue, alors deux primitives de f diffèrent d'une constante
\end{lm}

\begin{proof}
    Même preuve que précédemment.
\end{proof}

\begin{prop}
    Soit $\begin{array}{ccccc}
f : & \underset{Domaine}{D} & \to & \C
\end{array}$ continue, alors $f$ admet des primitives sur $D$ \ssi $\displaystyle{\int_\lambda f(z)\,dz=0}$ pour tout circuit $\lambda$ tracé dans $D$.
\end{prop}

\begin{proof}
    Voir p.25
\end{proof}

\subsection{Notion d'indice}

\begin{df}
    Etant donné un circuit $\lambda$ et un point $a\notin \text{supp}(\lambda)$.\\
    On appelle indice de $\lambda$ par rapport à $a$ : $$\text{Ind}(\lambda,a)=\frac{1}{2\pi i}\int_\lambda\frac{dz}{z-a}$$
\end{df}

\begin{prop}
Soit $\lambda$ un lacet, $a\notin\text{supp}(\lambda)$ et $\Omega=\C\backslash\text{supp}(\lambda)$
    \begin{enumerate}
        \item Ind$(\lambda,a)$ est constant quand a décrit un domaine inclus dans $\Omega$
        \item $\text{Ind}(\lambda,a)$ est constant sur tout chemin dont le support ne rencontre pas supp$(\lambda)$
        \item $\displaystyle{\lim_{|a|\to+\infty}\text{Ind}(\lambda,a)=0}$
        \item Ind$(\lambda,a)\in\Z$
    \end{enumerate}
\end{prop}

\begin{proof}
    Voir p.26
\end{proof}

\subsection*{Applications aux domaines étoilés}

Soit $\Delta$ domaine étoilé par rapport à $z_0$

\begin{prop}
    Si $\Delta$ est étoilé, tout contour triangulaire $\theta$ tracé dans $\Delta$, l'intérieur de $\theta$ est inclus dans $\Delta$.
\end{prop}

\chapter{Analycité des fonctions holomorphes}

\section{Théorème de Goursat}

\begin{tm}[Théorème (Goursat)]
    Toute fonction holomorphe sur un ouvert étoilé $\Omega$ y admet des primitives.\\
    On a vu qu'admettre des primitives sur un ouvert étoilé est équivalent à vérifier $\displaystyle{\int_\lambda f(z)\,dz=0}$ pour tout circuit triangulaire fermé $\lambda$ tracé dans $\Omega$
\end{tm}

\begin{proof}
    Voirp.27
\end{proof}

\section{Application : Théorème de d'Alembert}

\begin{tm}[Théorème (d'Alembert)]
    Tout polynôme à coefficients complexes de degré > 0 a au moins une racine.
\end{tm}

\begin{proof}
    Voir p.29
\end{proof}

\section{Formule intégrale de Cauchy pour un disque}

Soit D un disque ouvert de centre $z_0$ et de rayon $R$ et soit $\begin{array}{ccccc}
F : & D & \to & \C
\end{array}$ une fonction holomorphe. Pour tout $0<r<R$, on considère : $\gamma_r(t)=z_0+re^{it}\,(0\le t\le 2\pi)$

\begin{tm}[Théorème de Cauchy]
    $\forall a$ tel que $|a-z_0|<r$, $\displaystyle{F(a)=\frac{1}{2\pi i}\int_{\gamma_r} \frac{F(z)}{z-a}\,dz}$
\end{tm}

\begin{proof}
    Voir p.30
\end{proof}

\section{Analycité des fonctions holomorphes}

\subsection{Théorème}

\begin{tm}
    Toute fonction holomorphe sur un ouvert $\Omega$ est analytique sur $\Omega$
\end{tm}

\begin{proof}
    Voir p.31
\end{proof}

\begin{re}
    $$F^{(n)}(z)=\frac{n!}{2 \pi i}\int_\gamma \frac{F(u)}{(u-z)^{n+1}}\,du\,,\forall z\in D(a,r)$$
\end{re}

\begin{tm}
    Si $\begin{array}{ccccc}
F : & D & \to & \C
\end{array}$ est holomorphe sur un disque ouvert $D$, sa série de Taylor au centre de D converge absolument en tout point de D
\end{tm}

\subsection{Propriétés des coefficients}

Soit $\Omega$ ouvert et $\begin{array}{ccccc}
F : & \Omega & \to & \C
\end{array}$ holomorphe. Soit $a\in \Omega$ et $r>0$ tel que $\overline{D(a,r)}\subseteq \Omega$. Soit $\gamma(t)=a+re^{it},(0\le t\le 2\pi)$.\\
Si $\displaystyle{\sum_{n=0}^{+\infty}a_n(z-a)^n}$ est la série de Taylor de $F$ en $a$, $a_n= \frac{F^{(n)}(a)}{n!}$
$$\text{D'après la remarque } \displaystyle{a_n=\frac{1}{2\pi i}\int_\gamma\frac{F(u)}{(u-a)^{n+1}}\,du}.$$
$$\text{Donc } \displaystyle{a_n=\frac{1}{2\pi i}\int_{0}^{2\pi}\frac{F(a+re^{it})ir}{r^{n+1}e^{(n+1)it}}}$$.
$$a_nr^n=\frac{1}{2\pi}\int_0^{2\pi}F(a+re^{it})e^{-nit}\,dt$$
$a_nr^n$ s'interprète comme un coefficient de Fourrier pour $\begin{array}{ccccc}
& \R & \to & \C \\
& t & \mapsto & D(a+re^{it})  
\end{array}$

\begin{ra}
    Si $f$ est $2\pi$-périodique, intégrale sur $[0,2\pi]$, les coefficients de Fourrier exponentiels de $f$ sont les nombres complexes : $\displaystyle{c_n(f)=\frac{1}{2\pi}\int_0^{2\pi}f(t)e^{-int}\,dt},n\in \Z$\\
    Les coefficients de Fourrier trigonométriques sont :\\
    $$a_0(f)=\frac{1}{2\pi}\int_0^{2\pi}f(t)\,dt,a_n(f)=\frac{1}{\pi}\int_0^{2\pi}f(t)\cos(nt)\,dt$$
    $$b_n(f)=\frac{1}{\pi}\int_0^{2\pi}f(t)\sin(nt)\,dt,(n\ge1)$$
    $$c_n=\frac12(a_n-ib_n)\text{ et }c_{-n}=\frac12(a_n+ib_n)$$
\end{ra}

\subsection*{Séries de Fourrier}

$\displaystyle{\sum_{n\in\Z}c_n(f)e^{int}=a_0(f)+\sum_{n\ge 1}(a_n(f)\cos(nt)+b_n(f)\sin(nt))}$

\begin{prop}[Inégalité de Cauchy] 
Posons $\displaystyle{M(r):=\sup_{|z-a|=r}|F(z)|}$.
    $$|a_n|\le \frac{M(r)}{r^n}$$
\end{prop}

\begin{proof}
    En utilisant les propriétés précédentes
\end{proof}

\section{Quelques applications}

\begin{tm}[Théorème de Liouville]
    Toute fonction entière bornée est constante
\end{tm}

\begin{proof}
    Voir p.33
\end{proof}

\begin{cor}
    Tout polynôme à coefficients complexes de degré supérieur à 0 a au moins un zéro complexe.
\end{cor}

\begin{proof}
    Voir p.33
\end{proof}

\begin{tm}[Théorème de Morera]
    Si $\funto{f}{\Omega}{\C}$ est continue et vérifie "$\int_\theta f(z)\,dz=0$ pour tout circuit triangulaire tracé dans $\Omega$". Alors $f$ est holomorphe sur $\Omega$.
\end{tm}

\begin{proof}
    Voir p.34
\end{proof}

\chapter{Zéros des fonctions holomorphes}

\section{Étude au voisinage d'un zéro}

\subsection{Factorisation}

Soit $\funto{F}{\underset{ouvert}{\Omega}}{\C}$ holomorphe. On appelle zéro de $F$ un élément $z_0\in\Omega$ telle que $F(z_0)=0$


\begin{lm}
    Soit $z_0\in\Omega$ zéro de $F$. Les conditions suivantes sont équivalentes :\\
    \begin{itemize}
        \item $F(z)=0$ pour tout $z$ dans un voisinnage de $z_0$
        \item $\forall k\ge 0,F^{(k)}(z_0)=0$
    \end{itemize}
\end{lm}

\begin{proof}
    Voir p.35
\end{proof}

\begin{df}
    On appelle zéro isolé de $F$ tout zéro $z_0$ de $F$ qui ne vérifie pas les conditions du lemme.
\end{df}

\begin{prop}
    Si $z_0$ est un zéro isolé de $F$, il existe un unique entier positif $p$ tel que $F(z)=(z-z_0)^pF_1(z)$ où $\funto{F_1}{\Omega}{\C}$ holomorphe et $F_1(z_0)\neq 0$
\end{prop}

\begin{proof}
    Voir p.35
\end{proof}

\subsection{Ordre de multiplicité}

L'entier $p$ introduit dans la proposition précédente s'appelle l'ordre de multiplicité de $z_0$ ou simplement ordre de $z_0$. Dans la preuve précédente, il est caractérisé par les conditions :\\
$$F^{(k)}=0\text{ pour tout }k<p \text{ et }F^{(p)}(z_0)\neq 0$$
     Si p=1, le zéro est dit simple\\
$\left.\begin{array}{ll}
 \text{Si }p=2\text{ le zéro est dit double}\\  
 \text{Si }p=3\text{ le zéro est dit triple}\\
 \vdots\end{array}\right\} \text{ le zéro est dit multiple}$

\begin{nt}[Convention]
    Si $F(z_0)\neq 0$ on dira que $z_0$ est un zéro d'ordre $0$
\end{nt}

\begin{lm}
    Soient $f$ et $g$ deux fonctions holomorphes en $z_0$. Si $é_0$ est un zéro d'ordre $k$ pour $f$ et d'ordre $l$ pour $g$, alors $z_0$ est un zéro d'ordre $k+l$ pour $fg$.
\end{lm}

\begin{proof}
    $\left.\begin{array}{ll}
        f(z)=(z-z_0)^kF_1(z)\\
        g(z)=(z-z_0)^lG_1(z)
    \end{array}\right\}$ doù $(fg)(z)=(z-z_0)^{k+l}F_1(z)G_1(z)$
\end{proof}

\section{Distribution des zéros d'une fonction holomorphe}

\subsection{Principe des zéros isolés}

\begin{lm}
    Si $z_0$ est un zéro isolé de $\funto{F}{\underset{ouvert}{\Omega}}{\C}$ holomorphe, il existe un voisinage de $z_0$ sur lequel $F$ n'a pas d'autre zéro que ce point.
\end{lm}

\begin{proof}
    $F(z)=(z-z_0)^pF_1(z)$ avec $F_1(z_0)\neq 0$, $F_1$ étant continue et non nulle.\\
    $\exists r>0,|z-z_0|<r\implies F_1(z)\neq 0$. $z_0$ est alors le seul zéro de $F$ dans $D(z_0,r)=0$
\end{proof}

\subsection{Principe de prolongement}

\begin{lm}
    Soit $F$ holomorphe sur un domaine $\Delta\subseteq\C$. Si $F$ est nulle en tout point d'un disque ouvert inclus dans $\Delta$, $F$ est nulle sur $\Delta$
\end{lm}

\begin{prop}[Conséquence du Lemme]
    Si $\funto{F,G}{\Delta}{\C}$ holomorphe coïncident sur un disque ouvert contenu dans $\Delta$, on a $F=G$ sur tout $\Delta$
\end{prop}

\end{document}
